\section{Detailed Derivations}
\label{app:derivations}

This appendix gives the derivations underlying the deterministic equivalents used in the main text. We keep the notation consistent with Section~\ref{eq:linear_model}: $X\in\mathbb{R}^{n\times d}$, $\widehat{\Sigma}=X^\top X/n$, $\Sigma=\sum_{k=1}^d\lambda_k \uhat_k\uhat_k^\top$, and $b_k=\uhat_k^\top\btrue$.

\subsection{Ridge Decomposition}

For the ridge estimator in equation~\eqref{eq:ridge_estimator},
\begin{align}
    \bhat
    &= (\widehat{\Sigma}+\lambda I)^{-1}\widehat{\Sigma}\btrue
    + \sigma(\widehat{\Sigma}+\lambda I)^{-1}\frac{1}{n}X^\top\epsilon ,
    \label{eq:app_bhat_decomp}\\
    \bhat-\btrue
    &= -\lambda(\widehat{\Sigma}+\lambda I)^{-1}\btrue
    + \sigma(\widehat{\Sigma}+\lambda I)^{-1}\frac{1}{n}X^\top\epsilon .
    \label{eq:app_delta_beta}
\end{align}
The first term is the shrinkage of the signal induced by ridge regularization. The second term propagates response noise through the inverse empirical covariance.

\subsection{Deterministic Equivalent and the Effective Ridge Parameter}

In the proportional limit $d/n\to\gamma$, the ridge resolvent of $\widehat{\Sigma}$ is approximated by a deterministic resolvent of $\Sigma$ with an effective regularization parameter $\kappa=\kappa(\lambda)$ satisfying
\begin{equation}
    \kappa-\lambda
    =
    \gamma\kappa\frac{1}{d}\sum_{j=1}^{d}\frac{\lambda_j}{\kappa+\lambda_j}.
    \label{eq:app_kappa_fixed_point}
\end{equation}
The derivative of $\kappa$ with respect to $\lambda$ determines the finite-sample noise factor. Rewriting equation~\eqref{eq:app_kappa_fixed_point} as
\begin{equation}
    \lambda =
    \kappa\left[
    1-\gamma\frac{1}{d}\sum_{j=1}^{d}\frac{\lambda_j}{\kappa+\lambda_j}
    \right],
\end{equation}
we obtain
\begin{align}
    \frac{d\lambda}{d\kappa}
    &=
    1-\gamma\frac{1}{d}\sum_{j=1}^{d}\frac{\lambda_j}{\kappa+\lambda_j}
    +\gamma\frac{\kappa}{d}\sum_{j=1}^{d}\frac{\lambda_j}{(\kappa+\lambda_j)^2} \\
    &=
    1-\gamma\frac{1}{d}\sum_{j=1}^{d}\frac{\lambda_j^2}{(\kappa+\lambda_j)^2}.
\end{align}
With the unnormalized degrees of freedom
\begin{equation}
    \df(\kappa)=\sum_{j=1}^{d}\frac{\lambda_j^2}{(\lambda_j+\kappa)^2},
\end{equation}
and $\gamma=d/n$, this gives
\begin{equation}
    \frac{d\lambda}{d\kappa}=1-\frac{\df(\kappa)}{n},
    \qquad
    \frac{d\kappa}{d\lambda}=\frac{n}{n-\df(\kappa)}.
    \label{eq:app_dkappa}
\end{equation}

\subsection{Per-Eigenmode Coefficient Error}

Projecting equation~\eqref{eq:app_delta_beta} onto $\uhat_k$ and squaring gives three terms: signal bias, noise variance, and a signal-noise cross term. The cross term vanishes after averaging over the response noise. Thus
\begin{align}
    \mathbb{E}\left[(\uhat_k^\top(\bhat-\btrue))^2\right]
    &=
    \lambda^2
    \mathbb{E}\left[
    {\btrue}^{\top}(\widehat{\Sigma}+\lambda I)^{-1}
    \uhat_k\uhat_k^\top
    (\widehat{\Sigma}+\lambda I)^{-1}\btrue
    \right]
    \nonumber\\
    &\quad+
    \frac{\sigma^2}{n}
    \mathbb{E}\left[
    \operatorname{Tr}\left(
    (\widehat{\Sigma}+\lambda I)^{-1}
    \uhat_k\uhat_k^\top
    (\widehat{\Sigma}+\lambda I)^{-1}
    \widehat{\Sigma}
    \right)\right].
    \label{eq:app_perpc_start}
\end{align}

For the noise term, use
\begin{equation}
    (\widehat{\Sigma}+\lambda I)^{-1}\widehat{\Sigma}
    (\widehat{\Sigma}+\lambda I)^{-1}
    =
    \frac{d}{d\lambda}\left[\lambda(\widehat{\Sigma}+\lambda I)^{-1}\right].
\end{equation}
The deterministic equivalent gives
\begin{equation}
    \mathbb{E}\left[
    \uhat_k^\top
    (\widehat{\Sigma}+\lambda I)^{-1}\widehat{\Sigma}
    (\widehat{\Sigma}+\lambda I)^{-1}
    \uhat_k
    \right]
    \asymp
    \frac{d\kappa}{d\lambda}
    \frac{\lambda_k}{(\lambda_k+\kappa)^2}.
\end{equation}
Using equation~\eqref{eq:app_dkappa}, the noise contribution becomes
\begin{equation}
    \frac{\sigma^2}{n-\df(\kappa)}
    \frac{\lambda_k}{(\lambda_k+\kappa)^2}.
    \label{eq:app_noise_term}
\end{equation}

For the signal term, we use the two-resolvent deterministic equivalent. For deterministic matrices $A,B$ with bounded rank or bounded trace norm in the population eigenbasis,
\begin{align}
    &\operatorname{Tr}\left[
    A(\widehat{\Sigma}+\lambda I)^{-1}
    B(\widehat{\Sigma}+\lambda I)^{-1}
    \right]
    \nonumber\\
    &\quad\asymp
    \frac{\kappa^2}{\lambda^2}
    \operatorname{Tr}\left[
    A(\Sigma+\kappa I)^{-1}
    B(\Sigma+\kappa I)^{-1}
    \right]
    \nonumber\\
    &\qquad+
    \frac{\kappa^2}{\lambda^2}
    \frac{
    \operatorname{Tr}\left[
    A(\Sigma+\kappa I)^{-2}\Sigma
    \right]
    \operatorname{Tr}\left[
    B(\Sigma+\kappa I)^{-2}\Sigma
    \right]}
    {n-\df(\kappa)} .
    \label{eq:app_two_resolvent}
\end{align}
Taking $A=\btrue{\btrue}^{\top}$ and $B=\uhat_k\uhat_k^\top$ yields
\begin{align}
    &\lambda^2
    \mathbb{E}\left[
    {\btrue}^{\top}(\widehat{\Sigma}+\lambda I)^{-1}
    \uhat_k\uhat_k^\top
    (\widehat{\Sigma}+\lambda I)^{-1}\btrue
    \right]
    \nonumber\\
    &\quad\asymp
    \frac{\kappa^2}{(\lambda_k+\kappa)^2}b_k^2
    +
    \frac{\kappa^2}{n-\df(\kappa)}
    \frac{\lambda_k}{(\lambda_k+\kappa)^2}
    \sum_{j=1}^{d}
    \frac{\lambda_j b_j^2}{(\lambda_j+\kappa)^2}.
    \label{eq:app_signal_term}
\end{align}
Defining
\begin{equation}
    C_{\mathrm{sig}}(\kappa)=
    \sum_{j=1}^{d}
    \frac{\lambda_j b_j^2}{(\lambda_j+\kappa)^2},
\end{equation}
and combining equations~\eqref{eq:app_noise_term} and~\eqref{eq:app_signal_term} gives the main-text formula:
\begin{align}
    \mathbb{E}\left[(\uhat_k^\top(\bhat-\btrue))^2\right]
    &\asymp
    \frac{\kappa^2}{(\lambda_k+\kappa)^2}b_k^2
    +
    \frac{1}{n-\df(\kappa)}
    \frac{\lambda_k\kappa^2}{(\lambda_k+\kappa)^2}
    C_{\mathrm{sig}}(\kappa)
    \nonumber\\
    &\quad+
    \frac{\sigma^2}{n-\df(\kappa)}
    \frac{\lambda_k}{(\lambda_k+\kappa)^2}.
    \label{eq:app_perpc_final}
\end{align}

\subsection{Generalization Error}

For a fresh test input $\x\sim\mathcal{N}(0,\Sigma)$,
\begin{align}
    E_{\mathrm{gen}}
    &=
    \mathbb{E}\left[
    ((\bhat-\btrue)^\top\x)^2
    \right] \\
    &=
    \mathbb{E}\left[
    (\bhat-\btrue)^\top\Sigma(\bhat-\btrue)
    \right] \\
    &=
    \sum_{k=1}^{d}
    \lambda_k
    \mathbb{E}\left[
    (\uhat_k^\top(\bhat-\btrue))^2
    \right].
    \label{eq:app_gen}
\end{align}
Thus average prediction error weights coefficient error along each population PC by the natural-image variance $\lambda_k$.

\subsection{Accentuation Error and Alignment Ratio}

Accentuation uses the learned direction as a control path. Consider centered accentuation samples $\x=\alpha\bhat$, with $\mathbb{E}\alpha=0$. We choose the scale so that the model-predicted variance along the accentuation path matches the natural response variance:
\begin{align}
    \operatorname{Var}(\bhat^\top(\alpha\bhat))
    &=
    (\bhat^\top\bhat)^2\operatorname{Var}(\alpha)
    =
    {\btrue}^{\top}\Sigma\btrue,\\
    \operatorname{Var}(\alpha)
    &=
    \frac{{\btrue}^{\top}\Sigma\btrue}{(\bhat^\top\bhat)^2}.
\end{align}
The accentuation error against the true response is then
\begin{align}
    \Eacc
    &=
    \mathbb{E}_{X,\epsilon,\alpha}
    \left[
    \left(\alpha\bhat^\top(\bhat-\btrue)\right)^2
    \right] \\
    &=
    {\btrue}^{\top}\Sigma\btrue\cdot
    \mathbb{E}_{X,\epsilon}
    \left[
    \frac{(\bhat^\top\bhat-\bhat^\top\btrue)^2}
    {(\bhat^\top\bhat)^2}
    \right] \\
    &=
    {\btrue}^{\top}\Sigma\btrue\cdot
    \mathbb{E}_{X,\epsilon}
    \left[
    \left(1-
    \frac{\bhat^\top\btrue}{\bhat^\top\bhat}
    \right)^2
    \right].
    \label{eq:app_acc_exact}
\end{align}
This motivates the alignment ratio
\begin{equation}
    R=\frac{\bhat^\top\btrue}{\bhat^\top\bhat}.
\end{equation}

\subsection{Deterministic Equivalent for Mean Alignment}

The numerator concentrates around
\begin{align}
    N_{\det}
    &=
    \mathbb{E}[\bhat^\top\btrue]
    \asymp
    {\btrue}^{\top}(\Sigma+\kappa I)^{-1}\Sigma\btrue \\
    &=
    \sum_{k=1}^{d}
    \frac{\lambda_k}{\lambda_k+\kappa}b_k^2.
    \label{eq:app_Ndet}
\end{align}
The denominator must be computed as $\mathbb{E}[\bhat^\top\bhat]$, not as the squared deterministic mean of $\bhat$. Using equation~\eqref{eq:app_bhat_decomp} and the per-PC second moment,
\begin{align}
    D_{\det}
    &=
    \mathbb{E}[\bhat^\top\bhat]\\
    &\asymp
    \sum_{k=1}^{d}
    \left(\frac{\lambda_k}{\lambda_k+\kappa}\right)^2 b_k^2
    +
    \frac{\kappa^2 C_{\mathrm{sig}}(\kappa)+\sigma^2}{n-\df(\kappa)}
    \sum_{k=1}^{d}
    \frac{\lambda_k}{(\lambda_k+\kappa)^2}.
    \label{eq:app_Ddet}
\end{align}
The second term in equation~\eqref{eq:app_Ddet} is crucial: it contains both finite-sample signal variance and the corrected effective-sample-size denominator for the noise term. The leading deterministic alignment is
\begin{equation}
    R_{\det}=\frac{N_{\det}}{D_{\det}},
    \qquad
    \Eacc^{(0)}
    =
    {\btrue}^{\top}\Sigma\btrue(1-R_{\det})^2.
    \label{eq:app_Rdet}
\end{equation}

\subsection{Second-Order Correction: Variance of the Alignment Ratio}

The approximation in equation~\eqref{eq:app_Rdet} captures the squared bias term in
\begin{equation}
    \mathbb{E}[(1-R)^2]=(1-\mathbb{E}R)^2+\operatorname{Var}(R),
\end{equation}
but omits $\operatorname{Var}(R)$. When $R\approx 1$, this variance term can dominate the full accentuation error.

Let $N=\bhat^\top\btrue$, $D=\bhat^\top\bhat$, and $R_{\det}=N_{\det}/D_{\det}$. A delta-method expansion gives
\begin{equation}
    R-R_{\det}
    \approx
    \frac{(N-N_{\det})-R_{\det}(D-D_{\det})}{D_{\det}}
    =
    \frac{\delta Q}{D_{\det}},
\end{equation}
where $Q=N-R_{\det}D=\bhat^\top(\btrue-R_{\det}\bhat)$. Thus
\begin{equation}
    \operatorname{Var}(R)
    \asymp
    \frac{\operatorname{Var}(Q)}{D_{\det}^2}.
    \label{eq:app_varR_Q}
\end{equation}

Decompose $\bhat=g+h$, with
\begin{align}
    g &\asymp (\Sigma+\kappa I)^{-1}\Sigma\btrue,\\
    h &\sim
    \mathcal{N}\left(0,\frac{\sigma^2}{n}H\right),
    \qquad
    H=(\Sigma+\kappa I)^{-1}\Sigma(\Sigma+\kappa I)^{-1}.
\end{align}
Keeping the noise-driven fluctuations and centering the quadratic form,
\begin{equation}
    Q
    =
    h^\top(\btrue-2R_{\det}g)
    -
    R_{\det}\left(h^\top h-\mathbb{E}[h^\top h]\right).
\end{equation}
The first term is odd in $h$ and the second is even, so they are uncorrelated. Therefore,
\begin{align}
    \operatorname{Var}(Q)
    &=
    \frac{\sigma^2}{n}
    (\btrue-2R_{\det}g)^\top H(\btrue-2R_{\det}g)
    +
    2R_{\det}^2
    \left(\frac{\sigma^2}{n}\right)^2
    \operatorname{Tr}(H^2).
\end{align}
In the eigenbasis of $\Sigma$, this becomes
\begin{align}
    \operatorname{Var}(R)
    &\asymp
    \frac{1}{D_{\det}^2}
    \left[
    \frac{\sigma^2}{n}
    \sum_{k=1}^{d}
    \frac{\lambda_k b_k^2}{(\lambda_k+\kappa)^2}
    \left(1-\frac{2R_{\det}\lambda_k}{\lambda_k+\kappa}\right)^2
    +
    \frac{2R_{\det}^2\sigma^4}{n^2}
    \sum_{k=1}^{d}
    \frac{\lambda_k^2}{(\lambda_k+\kappa)^4}
    \right].
    \label{eq:app_varR_final}
\end{align}
The full second-order accentuation approximation is
\begin{equation}
    \mathbb{E}[\Eacc]
    \asymp
    {\btrue}^{\top}\Sigma\btrue
    \left[
    (1-R_{\det})^2+\operatorname{Var}(R)
    \right].
    \label{eq:app_acc_full}
\end{equation}
Equation~\eqref{eq:app_varR_final} captures the response-noise contribution to alignment variance. A complete finite-$n$ theory should also include random-design fluctuations of $\widehat{\Sigma}$, which empirically dominate at very low response noise and require a higher-order CLT for anisotropic Wishart quadratic forms.

\subsection{Why Low-Variance Modes Are Dangerous for Control}

For ordinary least squares, the population coefficient along PC $k$ can be written heuristically as
\begin{equation}
    w_{\mathrm{OLS},k}
    =
    \sigma_y\frac{1}{\sqrt{\lambda_k}}
    \operatorname{corr}(\uhat_k^\top X,y).
\end{equation}
For ridge regression,
\begin{equation}
    w_{\mathrm{ridge},k}
    =
    \sigma_y
    \frac{\sqrt{\lambda_k}}{\lambda_k+\lambda}
    \operatorname{corr}(\uhat_k^\top X,y).
\end{equation}
OLS therefore has explicit $1/\sqrt{\lambda_k}$ amplification in low-variance directions, while ridge suppresses the most extreme blow-up but still amplifies modes near $\lambda_k\approx\lambda$ or, in the finite-sample deterministic equivalent, near $\lambda_k\approx\kappa$. This is the linear mechanism behind high-frequency, off-manifold feature accentuations: low-variance directions barely affect prediction on natural images but can dominate the model gradient used for control.
